% Folder 136, batch 12 (archivists' pages 221-240).
% Pass: Opus 5.5 (claude-opus-5-5), 2026-10-03 — first pass, unchecked against the pages by a human.
%
% Pages 223, 227, 229, 231, 233, 235, 237 and 239 are unrelated versos (a
% computer listing and typed pages of another text, with nothing in his hand)
% and carry no \page{}. Pages 224-225 are loose sheets on perforated listing
% paper, in a paler ink, with their own proposition numbering.
\documentclass[11pt,a4paper]{article}
\input{../preamble/grothendieck}

\folder{136}
\batch{12}
\pages{221}{240}
\foldertitle{Complexe de De Rham à puissance divisée [conférence de 1976 à l’IHÉS] : notes manuscrites (s.d.).}
\dating{[à partir de 1975-1976]}
\watermark{Édition de démonstration}

\begin{document}

\page{221}
\note{la page s'ouvre au milieu d'une phrase, commencée avant ce lot.}
linéaire sur $\mathcal{M}$, on trouve une unique antidérivation $\delta$ de
$\dot\Lambda\mathcal{M}$ qui sur $\Lambda^1\mathcal{M}$ se réduise à $\varphi$, et on
a $\delta^2 = 0$. (N.B. On a $H_0(\dot\Lambda\mathcal{M}, \delta) = S/\varphi(\mathcal{M})$, et
si $\mathcal{M}$ est libre de t.f., alors $H_i(\dot\Lambda\mathcal{M}, \delta) = 0$ pour $i \neq 0$
ssi $\mathcal{M} \to S$ est « transversale » -- il suffit p.ex. que une
base de $\mathcal{M}$ définisse une suite $S$-régulière,
et c'est aussi néc. si $S$ noeth et $\varphi(\mathcal{M}) \subset \operatorname{rad} S$.)
Ceci dit, on obtient les complexes (2) et (4)
en prenant respectivement $S = \mathrm{Sym}^*M$ ou $S = \Gamma^*M$,
$\mathcal{M} = M \otimes_A S$, et en prenant l'hom $\mathcal{M} \to S$
\struck{\ill{}} déduit par extension des scalaires de l'hom can
$M \xrightarrow{\sim}$ \struck{$\mathrm{Sym}^1M$} $S^1 \subset S$
\struck{\ill{}} ($S^1 = \mathrm{Sym}^1M$ resp. $\Gamma^1M$).

Les opérateurs $d\delta$, $\delta d$ dans $\mathrm{Sym}^*\otimes\dot\Lambda$ resp.
$\Gamma^*\otimes\dot\Lambda$ ne sont pas bien jolis, mais on voit de
suite (tout anticommutateur de antidérivations étant
une dérivation) que
\[ \theta = d\delta + \delta d \]
est une dérivation, \struck{et}; en changeant \struck{\ill{}} le
\uncertain{bidegré} \struck{\ill{} $S^* = \Gamma^*$ et compatible aux \ill{}}, et qui sur
$S^1 = \mathrm{Sym}^1$ ou $\Gamma^1$ \add{(de plus si $S^* = \Gamma^*$, il est compatible aux puiss.\ div.\ sur $\Gamma^*$)}
et sur $\dot\Lambda^1$ est l'identité; (Donc on en conclut facilement)
que l'on a
\[ (d\delta + \delta d)x = \theta x = nx \qquad \text{$x$ de degré total $n$,} \]
\[ \text{i.e. } x \in \textstyle\sum_{p+q=n} S^p \otimes \Lambda^q . \]
Cela montre donc que si $A$ est de car.\ zéro (ou même que $n$
inversible $\forall n > 0$), alors $H^n_{\mathrm{DR}}(S^*\otimes\dot\Lambda)$ et $H^n_{\mathrm{Kosz}}(S^*\otimes\dot\Lambda)$
sont nuls quel que soit $n > 0$
($S^*$ étant $\mathrm{Sym}^* \simeq \Gamma^*$) (sans hyp.\ de platitude!). Plus gén., $H^n_{\mathrm{DR}}$ et $H^n_{\mathrm{Kosz}}$
sont nuls pour \struck{tout} $n$ tel que $n.1_A$ inv.\ dans $A$.
\note{la lecture des lignes biffées et ajoutées en interligne autour de « en changeant … bidegré » est très incertaine ; l'ordre des fragments est reconstitué.}

\page{222}
Revenant au cas $A$ de car.\ $0$, que peut-on dire de
$H^0_{\mathrm{DR}}$ et $H_0^{\mathrm{Kosz}}$? Pour $H_0$, on sait déjà que c'est $\simeq A$;
est-il vrai que l'hom évident $A \to H^0_{\mathrm{DR}}$ est aussi
un isom?, \struck{\ill{}} i.e.\ que $\forall f \in \mathrm{Sym}^*(M)$ tel que
$df = 0$ soit \struck{\ill{}} une constante, ?, i.e.\ \add{($f \in \mathrm{Sym}^k(M)$, si $k>0$)}
$k > 0$ et $df = 0 \Rightarrow f = 0$, i.e.
\[ \mathrm{Sym}^kM \xrightarrow{\ d\ } (\mathrm{Sym}^{k-1}M) \otimes M \]
\add{($M$ étant ici $\simeq \dot\Lambda^1 M$)} est injectif? Or composant avec
$\mathrm{Sym}^{k-1}(M) \otimes M \xrightarrow{\ \delta\ } \mathrm{Sym}^kM$ \add{($M \simeq \mathrm{Sym}^1M$)}
on trouve la multiplication par $k$ (car $\theta = d\delta + \delta d$ est égal
à $\delta d$ sur $\mathrm{Sym}^kM$ puisque $\delta$ y est nulle), donc inj, ok.
(Plus gén., inj.\ pourvu que $k$ inv.\ dans $A$.)

\emph{Dualité.} Supposons $M$ et $M'$ deux $A$-modules accouplés
vers $A$,
\[ M \otimes M' \to A ; \]
alors on en conclut des accouplements (degré par degré)
\[ \mathrm{Sym}^*M \times \Gamma^*M' \to A, \qquad \dot\Lambda M \times \dot\Lambda M' \to A \]
qui est une dualité parfaite si $M$ est projectif t.f.\ et
$M'$ s'identifie à son dual. On en conclut des accouplements
\[ \text{\struck{$\mathrm{Sym}$}}\ S^*\Lambda^*(M) \times \Gamma^*\Lambda^*(M') \to A . \]
Ceci posé, on vérifie que les op.\ $d_M$ et $\delta_{M'}$, $\delta_M$ et $d_{M'}$
sont \struck{\ill{}} en transposition par ces accouplements.
D'ailleurs, on vérifie aussi que l'op.\ de multipli\-cation
d'un côté est en dualité avec l'op.\
de comultiplication de l'autre. Donc dans le cas
de dualité parfaite, les opérations $d_M$ et $\delta_{M'}$, $\delta_M$ et $d_{M'}$,
multiplication et comultiplication, se déterminent mutuellement.

\page{224}
\marginal{au crayon, en travers de la marge : Feuille volante \uncertain{paumée}…}
\note{feuillet isolé, d'une encre plus pâle, qui ne fait pas suite à la p. 222 ; il porte sa propre numérotation (« Prop 6 », renvois aux « prop.\ 3, 5 » et « prop 4 » d'un texte qui n'est pas dans ce lot).}
\emph{Prop 6.} Soit $\mathfrak{B}$ une catégorie \ill{}.
Alors le foncteur
\[ ND_\bullet : \mathfrak{B}^*_* \to \mathfrak{B}_\bullet \]
est donné (à isom.\ can.\ près) par
\[ ND_i(L_*) \simeq \mathrm{Hom}_{k,*}(\underbrace{\dot\Lambda^i\mathfrak{L}_*}_{\mathfrak{L}^i_*}, L_*) \]
(pour la définition de $\mathrm{Hom}_{k,*}$, cf.\ prop.\ 3, 5),
\struck{\ill{}}
\[ ND_\bullet(L_*) \simeq \mathrm{Hom}_{k,*}(\underbrace{\dot\Lambda\mathfrak{L}_*}_{\mathfrak{L}^\bullet_*}, L_*) ; \]
l'opérateur diff.\ sur $ND_\bullet(L_*)$ provient de
l'op.\ différentiel sur $T_*$ et sur $\dot\Lambda\mathfrak{L}_*$.

En effet, \struck{\ill{}} posant $L_\bullet = ND_\bullet(L_*)$,
\add{et \struck{\ill{}}} $\mathfrak{L}(i)_\bullet = ND_\bullet(\dot\Lambda^i\mathfrak{L}_*)$, alors on a par
\uncertain{la symétrie} de Dold-Puppe
\[ \mathrm{Hom}_{k,*}(\dot\Lambda^i\mathfrak{L}_*, L_*) \simeq \mathrm{Hom}_{k_\bullet}(\mathfrak{L}(i)_\bullet, L_\bullet) . \]
Or $\mathfrak{L}(i)_\bullet \simeq C_\bullet^{i-1}$ pour $i \geq 1$ (prop 4), et
$\mathfrak{L}(0)_\bullet = k$ placé en degré $0$;
on a bien
\[ \begin{cases} \mathrm{Hom}_{k_\bullet}(C_\bullet^{i-1}, L_\bullet) \simeq L_i & (i \geq 1) \\
\mathrm{Hom}_{k_\bullet}(k[0]_\bullet, L_\bullet) \simeq L_0 \end{cases} \]
et il reste à faire de petites vérifications
de compatibilité…

\emph{Cor.} Le foncteur $ND^\bullet : \mathfrak{B}^* \to \mathfrak{B}^\bullet$ s'explicite
comme
\[ ND^\bullet(L^*) \simeq \dot\Lambda\mathfrak{L}_* \otimes_{k,*} L^* \]
[où pour $M_* \in \mathrm{Ab}_{k_*}$ et $L^* \in \mathfrak{B}^*$,
$M_* \otimes_{k_*} L^*$ désigne l'objet de $\underline{\mathrm{Hom}}(\mathfrak{B}, (\mathrm{Ab}))$ \note{indice de $\mathrm{Ab}$ illisible}
\[ b \longmapsto \mathrm{Hom}_{\mathrm{Ab}_{k_*}}(M_*, \mathrm{Hom}_{\mathfrak{B}}(L^*, b)) \ ] \]
Il suffit d'appliquer la prop.\ 6 à $\mathfrak{B}^\circ$.

\page{225}
\note{feuillet du même type que la p.\ 224, qui semble le précéder dans l'ordre de rédaction (la « Prop 6 » y est ébauchée puis biffée) ; il commence au milieu d'un énoncé.}
\marginal{en haut à gauche, en oblique : \ill{} pour un foncteur \ill{} sur $\mathfrak{B}$ \ill{} Kan-\ill{}}
\[ F(x) \simeq \mathrm{Hom}^*_k(\alpha^*(x), C^*_F), \qquad C^*_F \overset{\mathrm{dfn}}{=} F(C^*) \]
\note{après la formule, un ou deux signes biffés.}
où la définition des $\mathrm{Hom}^*$ se fait comme
dans la prop.\ 3.

\note{le passage suivant, jusqu'à la flèche $F(C^i) = C^i \to \dots$, est entouré d'un trait et barré de traits obliques.}
\struck{$F(x) \otimes \alpha^{[n]}(x) \to C^{[n]}_F$, avec $\simeq \mathrm{Hom}(x, C^{[n]})$ et $C^{[n]}_F = F(C^{[n]})$.
On procède comme pour la prop.\ 3 pour
définir la flèche dans $\widehat{\mathfrak{B}}$, et pour
montrer que c'est un iso, on se ramène
au cas où $x$ \uncertain{est de la} forme $C^i_\xi$, $C$ fixe,
prenons \ill{}
$F(C^i) = C^i \to \mathrm{Hom}^*_k($\ill{}$C^*)$.}
\marginal{à gauche du passage biffé : par définitions des $\alpha^*$ et \struck{\ill{}} DP adjoint, \struck{\ill{}} et commutation des DP aux foncteurs \uncertain{etc.}}

En fait, on trouve tout de suite que si
$K^\bullet \in \mathrm{Ab}_k^\bullet$, $C^\bullet \in \mathfrak{B}^\bullet$, $K^* = DP(K^\bullet)$ et $C^* = DP(K^*)$\note{sic, pour $DP(C^\bullet)$ sans doute},
alors on a un iso.\ can.\ bifonctoriel dans $\widehat{\mathfrak{B}}$
\[ DP : \mathrm{Hom}^\bullet_k(K^\bullet, C^\bullet) \xrightarrow{\ \sim\ } \mathrm{Hom}^*_k(K^*, C^*), \]
et on conclut par prop 3.

\note{le passage suivant est encadré et barré de traits obliques.}
\struck{\emph{Prop 6} (*) Le foncteur
$ND : \mathrm{Ab}_{k_*} \to \mathrm{Ab}_{k_\bullet}$
est donné par
$(ND)(L_*)$ $\simeq \mathrm{Hom}_{k,*}(C^\bullet_*, L_*)$,
i.e.\ $ND(L_*)_i$ $\simeq \mathrm{Hom}_{k,*}(C^i_*, L_*)$ pour $i \geq 0$
($C^i_* = \dot\Lambda^i\mathfrak{L}_*$, $i \geq 1$),
l'hom \add{bord} $ND(L_*)_i$ $\to ND(L_*)_{i-1}$ étant
le transposé de $C^{i-1}_*$ $\to C^i_*$.}

\page{226}
\section{Quelques calculs d'homotopie}

Soit $k$ un anneau (de base \uncertain{commun}),
$\mathfrak{L}_* = \mathfrak{L}^k_*$ le $k$-module ½simplicial $I \mapsto k^I$,
$k_*$ le $k$-module ½simpl.\ constant défini par $k$, $I \mapsto k$.
On a une application diagonale
\[ k_* \to \mathfrak{L}_* \]
d'où un quotient
\[ \Psi_* = \mathfrak{L}_*/k_* : I \mapsto k^I/k . \]

\emph{Proposition 1.} $\mathfrak{L}_*$ est homotope à $0$ \struck{\ill{}} (\add{i.e.\ sur $k$} \uncertain{module} ½simplicial $0$).

Par le \uncertain{yoga} de Puppe-Dold, il suffit de voir que
le complexe de chaînes $\mathfrak{L}^!_*$ associé à $\mathfrak{L}$ est homot.\ à
$0$. Pour ceci, on se rappelle que
\[ \mathfrak{L}^!_n = \{ x \in \mathfrak{L}_n \mid \partial_1(x) = \dots = \partial_n(x) = 0 \} \]
d'où aussitôt
$\mathfrak{L}^!_n = \begin{cases} 0 & \text{si } n \geq 2 \\ k & \text{si } n = 0, 1 \end{cases}$,
l'op.\ \struck{\ill{}} bord
\[ \partial = \partial_0 : \mathfrak{L}^!_1 \to \mathfrak{L}^!_0 \qquad (\mathfrak{L}^!_1 \simeq k,\ \mathfrak{L}^!_0 \simeq k) \]
étant l'identité, d'où la conclusion.

\emph{Corollaire 1.1.} Pour tt foncteur $F : \mathrm{Mod}(k) \to \mathcal{C}$
($\mathcal{C}$ une catégorie \ill{}) $F(\mathfrak{L}_*)$ est homotope à $F(0)$.
En particulier, \add{$(\Gamma^i_k\mathfrak{L}_*, \mathrm{Sym}^i_k\mathfrak{L}_*)$}
\[ \dot\Lambda^i_k\mathfrak{L}_* \text{ est homotope à } \dot\Lambda^i_k(0) = \begin{cases} 0 & \text{si } i \neq 0 \\ k_* & \text{si } i = 0 \end{cases} \]
\add{(de même $\Gamma^i_k(0)$, $\mathrm{Sym}^i_k(0)$)}
\struck{$\Gamma^i_k\mathfrak{L}_*$ \ill{} $\Gamma^i_k(0) = \dots$}
\note{une seconde ligne de l'accolade, pour $\Gamma^i_k\mathfrak{L}_*$, est biffée d'une rature serrée ; l'énoncé est reporté dans l'ajout entre parenthèses.}

\emph{Corollaire 1.2.} Les groupes d'homotopie de $\Psi_* =
\mathfrak{L}_*/k_*$ sont
\[ \pi_i(\Psi_*) \simeq \pi_{i-1}(k_*) = \begin{cases} 0 & \text{si } i \neq 1 \\ k & \text{si } i = 1 \end{cases} \]
En d'autres termes, $\Psi_*$ \struck{\ill{}} est un $K_*(k, 1)$.
[En fait, on a $(\Psi^!_*) =$ complexe réduit à $k$ placé en degré $1$.]
On conclut de ceci \struck{et du \uncertain{yoga} de Dold-Puppe}

\emph{Corollaire 1.3.}
\[ \pi_j(\dot\Lambda^i(\mathfrak{L}_*/k_*)) = \begin{cases} 0 & \text{si } j \neq i \\ k & \text{si } j = i \end{cases} \]
(i.e.\ $\dot\Lambda^i\mathfrak{L}_*/k_*$ est un $K_*(k, i)$).

\page{228}
Dans le langage \uncertain{de Dold-Puppe} des foncteurs
dérivés de $\dot\Lambda$, compte tenu de $(\mathfrak{L}_*/k_*)^! = k[1]$,
l'assertion équivaut à
\[ (L\dot\Lambda^i)(k[1]) \simeq \dot\Lambda^i k[1] = k[\text{-}i] \]
\note{la fin de la formule se lit $k[-i]$ ou $k[i]$ ; le N.B. plus bas écrit $k[i]$.}
-- plus généralement, si on se donne un
complexe de chaînes de $k$-modules $L_\bullet$
projectifs en chaque degré,
\marginal{est-ce vrai?}
\[ (L\dot\Lambda)(L_\bullet) \xrightarrow{\ \sim\ } \dot\Lambda(L_\bullet) . \]
On peut aussi noter que la filtration de $\mathfrak{L}_*$ par
\[ 0 \to k_* \to \mathfrak{L}_* \to \Psi_* \to 0 \qquad (\Psi_* \overset{\mathrm{dfn}}{=} \mathfrak{L}_*/k_*) \]
définit pour les $i \geq 0$ une filtration
\[ 0 \to \underbrace{\dot\Lambda^{i-1}\Psi_* \otimes k_*}_{\dot\Lambda^{i-1}\Psi_*} \to \underbrace{\dot\Lambda^i\mathfrak{L}_*}_{\text{flasque}} \to \dot\Lambda^i\Psi_* \to 0 \]
d'où
\[ 0 \to (\dot\Lambda^{i-1}\Psi_*)^! \to \underbrace{(\dot\Lambda^i\mathfrak{L}_*)^!}_{\text{flasque}} \to \dot\Lambda^i\Psi_* \to 0 \]
\note{cette dernière suite est entourée d'un trait.}
d'où la conclusion, car on a aussi $\pi_j(\dot\Lambda^i\Psi_*) \simeq \pi_{j-1}(\dot\Lambda^{i-1}\Psi_*)$
par suite exacte d'homotopie.

\textbf{NB}. Il est probable que $(\dot\Lambda^i\Psi_*)^!$ est \struck{\ill{}} $k[i]$
($k$ placé en dim $i$). On n'en aura pas besoin.

b) Les mêmes arguments montrent que
$\Gamma^i(\Psi^*)$ et \struck{\ill{}} $\mathrm{Sym}^i(\Psi^*)$ sont des $K(k, i)$\note{une lettre biffée précède aussi $\Psi^*$ dans $\Gamma^i(\Psi^*)$.}
\add{\ill{} $\otimes \Psi^*$} (mais nettement moins « économiques »!)

\emph{Cor.\ 1.4.} Si $A_*$ est un $k$-module ½simplicial, \ill{}
\[ \pi_\ell(A_* \otimes \dot\Lambda^i\Psi_*) \simeq \pi_{\ell-i}(A_*) \]
[de façon plus précise, \add{on doit avoir} $A_* \otimes \dot\Lambda^i\Psi_* \simeq$ \struck{$\Omega^i$} $\mathrm{Susp}^iA_*$,
où Susp est la « suspension additive »]

\page{230}
Par Eilenberg-Zilber, l'\struck{\ill{}}homologie de
$C_*(A_* \otimes B_*)$ (\add{$A_*$, $B_*$} deux $k$-modules ½simpliciaux \struck{\ill{}})
\add{$C_*(A_* \otimes B_*) = A_* \otimes B_*$}
\struck{complexe de chaînes} \add{op.} bord $\partial = (\sum\partial_n)$, $\partial_n = \sum_0^n (-1)^i\partial_{i,n}$, --
c'est donc, par Dold-Puppe, aussi $\pi_*(A_* \otimes B_*)$,
se calcule en prenant \add{l'homologie} \struck{du} produit tensoriel
\struck{des complexes} $C_*(A_*) \otimes C_*(B_*)$, qui est
lui-même « \add{chaîne-}\struck{homotope} », par Dold-Puppe,
à \struck{\ill{}} $A^!_* \otimes B^!_*$ (car $A^!_* \xrightarrow{\ \sim\ } A_*$ \add{chaîne-homotopisme},
$B^!_* \to B_*$); \add{a son homologie} \struck{laquelle} qui se calcule par
Künneth. Ici $B^!_* = (\dot\Lambda^i\Psi_*)^! \approx k[i]$, donc
on \struck{doit avoir} a
\[ \underbrace{H_\ell(A^!_* \otimes B^!_*)}_{\pi_\ell(A_* \otimes B_*)} \simeq \underbrace{H_{\ell-i}(A^!_*)}_{\pi_{\ell-i}(A_* \otimes B_*)} \]
\note{sous $H_{\ell-i}(A^!_*)$ on lit $\pi_{\ell-i}(A_* \otimes B_*)$, sans doute pour $\pi_{\ell-i}(A_*)$.}

\emph{Cor.\ 1.5.}
Considérons p.ex.\ le complexe \add{½simplicial} \struck{de chaînes} De Rham
à puissances divisées
\[ A^{**}_* \simeq \Gamma^*(\mathfrak{L}_*) \otimes_k \dot\Lambda\Psi_* \qquad \text{(degré total, degré ext)} \]
on trouve alors
\[ \pi_\ell(A^{*i}_*) \simeq \pi_{\ell-i}(\underbrace{A^{*0}_*}_{\Gamma^*(\mathfrak{L}_*)}) \rightleftarrows \pi_{\ell-i}(\underbrace{A^{00}_*}_{\Gamma^*(0) = k_*}) = \begin{cases} 0 & \text{si } \ell \neq i \\ k & \text{si } \ell = i \end{cases} \]
En particulier \struck{\ill{}}
\[ \pi_\ell(\underbrace{A^{p,i}_*}_{\Gamma^{p-i}_k\mathfrak{L}_* \otimes \dot\Lambda^i\Psi_*}) \simeq \pi_{\ell-i}(\underbrace{A^{p-i,0}_*}_{\Gamma^{p-i}_*\mathfrak{L}_*}) \simeq \begin{cases} 0 & \text{si } \ell \neq i \text{ ou } p \neq i \\ k & \text{si } i = \ell = p \end{cases} \]
donc, pour $p$ fixé, $i$ variable, on trouve que les
\struck{comp.} \add{composantes}
$A^{p,0}_*, A^{p,1}_*, \dots, A^{p,p-1}_*$ (en degré $< p$) de $A^{p*}_*$
(qui est une \uncertain{résolution} de $k_*$) sont
flasques, la seule composante non flasque étant
$A^{p,p}_*$ (ayant comme seul gr.\ d'homot.\ non nul
$\pi_p(A^{p,p}_*) = k$). On peut donc considérer que

\page{232}
ce complexe est déduit d'une résolution
flasque $\widetilde{\mathfrak{L}}^{p*}_*$ de \struck{\ill{}} $k_*$ en la « tronquant »
en \struck{dimension} degré $\leq p$, i.e.\ en prenant $\mathfrak{L}^{p*}_* = \tau^{\leq p}\widetilde{\mathfrak{L}}^{p*}_*$,
\struck{\ill{}} i.e.
\[ \mathfrak{L}^{pi}_* = \begin{cases} \widetilde{\mathfrak{L}}^{p,i}_* & \text{si } i < p \\ Z^p(\widetilde{\mathfrak{L}}^{p,*}_*) & \text{si } i = p \\ 0 & \text{si } i > p \end{cases} \]
Par suite, pour tout $\mathcal{X}$ semi-simplicial, on
aura
\[ \underbrace{\mathrm{Hom}(\mathcal{X}, \mathfrak{L}^{p*}_*)}_{C^{p*}_{\mathrm{DR\,pd}}(\mathcal{X}, k)} \simeq \mathrm{Hom}(X, \tau^{\leq p}\widetilde{\mathfrak{L}}^{p*}_*) \simeq \tau^{\leq p}\Gamma(X, \widetilde{\mathfrak{L}}^{p*}_*) \]
\[ \simeq \tau^{\leq p}\mathbb{R}\Gamma(\mathcal{X}, k) \]
\note{sous $\mathrm{Hom}(\mathcal{X}, \mathfrak{L}^{p*}_*)$, un signe d'égalité marqué « dfn » ; la lettre surmontée d'un tilde se lit comme le $\mathfrak{L}$ des pages précédentes, sans certitude.}
en particulier
\[ H^{p*}_{\mathrm{DR\,pd}}(\mathcal{X}, k) \to H^*(\mathcal{X}, k) \]
\struck{est un isom} tel que
\[ H^{p,i}_{\mathrm{DR\,pd}}(\mathcal{X}, k) \begin{cases} \simeq H^i(\mathcal{X}, k) & \text{si } i \leq p \\ = 0 & \text{si } i > p . \end{cases} \]

Soit $S$ anneau avec idéal $S^+$ à p.d., muni d'un
\struck{système de générateurs} $t \in S^+$. Considérons \uncertain{le complété}
$S$ \ill{} pour la top.\ engendrée
par \uncertain{ces} idéaux
\[ S^{[i]} = \text{idéal engendré par les } t^{[j]} \ (j \geq i) \]
-- ou on suppose $S$ séparé et complet pour
la top.\ définie par une suite décroissante
d'idéaux $S^{[i]}$, chacun étant l'adhérence
de l'idéal engendré par $t^{[j]}$ ($j \geq i$).

\page{234}
On considère l'anneau ½simplicial
\[ \widehat{A}^*_*(S, \gamma^*, t) \overset{\mathrm{dfn}}{=} \widehat{A}^{**}_* \]
par
\[ \widehat{A}^*_n \simeq S\{\{X_0, \dots, X_n\}\}/(\text{idéal à p.d.\ engendré par } \textstyle\sum X_i - t) \otimes_S \dot\Lambda[dX_0, \dots, dX_n]/\text{idéal}(\textstyle\sum dX_i) \]
\[ (\simeq S\{\{Y_0, \dots, Y_n\}\}[dY_0, \dots, dY_n]) \]
\note{dans la parenthèse « idéal à p.d.\ … », quelques mots biffés avant « engendré par » ; la lettre lue $t$ ressemble à une majuscule $T$.}
On a donc
\[ \widehat{A}^*_* \simeq \widehat{A}^0_* \otimes_S \dot\Lambda\Psi^S_* \]
d'où
\[ \pi_\ell(\widehat{A}^i_*) \simeq \pi_\ell(\widehat{A}^0_* \otimes \dot\Lambda^i\Psi^S_*) \simeq \pi_{\ell-i}(\widehat{A}^0_*) \]
et il faut donc étudier $\pi_*(\widehat{A}^0_*)$.

On considère sur $\widehat{A}^0_*$ \struck{\ill{}} la filtration
décroissante par idéaux fermés stables par p.d.\
« engendrée » par la condition que \struck{\ill{}} pour tout $n$,
les $X_i$ et $t$ sont en filtration $\geq 1$. On a
\[ \mathrm{Fil}^0(\widehat{A}^0_n) \simeq \widehat{A}^0_n \]
\begin{align*}
\mathrm{Fil}^p(\widehat{A}^0_n) &\simeq \Bigl\{ \textstyle\sum a_{i_0 \dots i_n} X_0^{[i_0]} \cdots X_n^{[i_n]} \Bigm| a_{i_0 \dots i_n} \in S^{[p - \sum_0^n i_\alpha]} \Bigr\} \\
&= \Bigl\{ \textstyle\sum a_{i_0, \dots, i_{n-1}} Y_0^{[i_0]} \cdots Y_{n-1}^{[i_{n-1}]} \Bigm| a_{i_0 \dots i_{n-1}} \in S^{[p - \sum_0 i_\alpha]} \Bigr\}
\end{align*}
Bien sûr, la topologie qui provient de \ill{} sur $\widehat{A}_n$ est
aussi celle définie par cette filtration, et par
suite $\widehat{A}^0_n$ est séparé et complet pour cette
filtration. Un calcul immédiat donne,
posant
\[ k = S/S^{[1]} \]
\[ \mathrm{Gr}^p(\widehat{A}^0_*) \simeq \Gamma^p_k(\mathfrak{L}^k_*) \]
\add{$\simeq \widehat{A}^{p0}_*(\mathrm{Gr}^*(S), \widetilde{\gamma}^*, \widetilde{t})$, $k$-½simplicial gradué, gradué de $k[T]$}
\struck{\ill{}} si l'hom $k[T] \to \mathrm{Gr}^*(S)$ est un iso,
d'où, dans cette hypothèse \add{d'isomorphisme},
\[ \pi_i(\mathrm{Gr}^p(\widehat{A}^0_*)) = \begin{cases} 0 & \text{si } p \neq 0,\ i \neq 0 \\ k & \text{si } i, p = 0 \end{cases} \]
(donc $\mathrm{Gr}^p(\widehat{A}^0_*)$ flasque si $p > 0$)
d'où on conclut que \struck{\ill{}}
\struck{$\mathrm{Fil}^1\widehat{A}^0_* \simeq \widehat{A}^0_*$}
\note{la page s'arrête sur cette conclusion biffée ; la p.\ 236 reprend.}

\page{236}
\struck{On trouve donc
$H^i(C^*_{\mathrm{DR}}(\mathcal{X}, \widehat{A}^*_*)) = H^0(\mathcal{X}, S)$ si $i = 0$,
$H^i(\mathcal{X}, S^+)$ si $i \neq 0$}
\note{cette accolade est barrée de traits obliques.}
\[ \& \begin{cases} \mathrm{Fil}^p(\widehat{A}^0_*) \text{ flasque pour tt } p \geq 1 \\
\pi_i(\widehat{A}^0_*) = \begin{cases} 0 & \text{si } j \neq 0 \\ k & \text{si } i = 0 \end{cases} \quad \text{si } p = 0 \end{cases} \]

Considérons la filtration \add{½simpl.} \ill{} sur $\widehat{A}^*_*$
prolongeant celle de $\widehat{A}^0_*$ et donnant \struck{\ill{}}
à $\widehat{A}^1_*$ la filtration $1$, donc
\[ \mathrm{Fil}^p\widehat{A}^*_* = \sum_{0 \leq i \leq p} \mathrm{Fil}^{p-i}(\widehat{A}^0_*) \otimes_S \dot\Lambda^i\Psi_* \]
C'est stable par l'opérateur différentiel.
On a donc
\begin{align*}
\pi_\ell(\mathrm{Fil}^p(\widehat{A}^i_*)) &= \pi_\ell(\mathrm{Fil}^{p-i}(\widehat{A}^0_*) \otimes_S \dot\Lambda^i\Psi_*) \\
&= \pi_{\ell-i}(\mathrm{Fil}^{p-i}(\widehat{A}^0_*)) \\
&= \begin{cases} 0 & \text{si } p \leq i \text{ et } \ell \neq i \\ k & \text{si } p \leq i = \ell \end{cases}
\end{align*}
\note{la colonne de droite porte plusieurs essais biffés de ces conditions (« $\ell \neq p$ ou … », « $0$ si $\ell - i \neq$ … ») ; seule la dernière rédaction est donnée.}

Donc: \struck{on en conclut la suite spectrale (pour \ill{} fixé)}
si on regarde $\mathrm{Fil}^p(\widehat{A}^*_*)$ comme
\struck{$E_2^{rs} = H^r$ \ill{} $(\mathcal{X}$, \ill{}}
une résolution ½simpliciale de $S^{[p]}_* = \mathrm{Fil}^pS$,
on voit que tous les termes $\mathrm{Fil}^p(A^i_*)$
de cette résolution \add{en degré $< p$} sont flasques, \struck{\ill{}}
\struck{\ill{}} ($j \geq p$) $\mathrm{Fil}^p(A^p_*) \simeq \widehat{A}^0_* \otimes_S \dot\Lambda^p\Psi_*$ dont
le seul gr.\ d'homotopie non nul est
$\pi_p(\cdot) \simeq k$. On en conclut

\begin{tikzcd}[column sep=small]
H^i(\underbrace{\mathrm{Fil}^pC^*_{\mathrm{DR}}(\mathcal{X}, \widehat{A}^*_*)}_{C^*_{\mathrm{DR}}(\mathcal{X}, \mathrm{Fil}\widehat{A}^*_*)}) \arrow[r, "\simeq"] \arrow[d] & H^i(\mathcal{X}, \underbrace{S^{[p]}}_{\mathrm{Fil}^pS}) \arrow[d] \\
H^i(C^*_{\mathrm{DR}}(\mathcal{X}, \widehat{A}^*_*)) \arrow[r] & H^i(\mathcal{X}, S)
\end{tikzcd}

si $i \leq p$.
\note{sous le premier terme, un signe « $\simeq$ dfn » relie $\mathrm{Fil}^pC^*_{\mathrm{DR}}(\mathcal{X}, \widehat{A}^*_*)$ à $C^*_{\mathrm{DR}}(\mathcal{X}, \mathrm{Fil}\,\widehat{A}^*_*)$ ; la flèche horizontale du haut est un isomorphisme pour $i \leq p$.}

\page{238}
\uncertain{qui} s'écrit aussi

\begin{tikzcd}[column sep=small]
H^i(\mathrm{Fil}^pC^*_{\mathrm{DR}}(\mathcal{X}; S, t, \gamma^*)) \arrow[r, "\sim"] \arrow[d] & H^i(\mathcal{X}, \mathrm{Fil}^pS) \arrow[d] \\
H^i(\mathrm{Fil}^qC^*_{\mathrm{DR}}(\mathcal{X}; S, t, \gamma^*)) \arrow[r] & H^i(\mathcal{X}, \mathrm{Fil}^qS)
\end{tikzcd}

\marginal{entouré, en haut à droite : iso si $i \leq p$ ; à droite de la seconde ligne : iso si $i \leq q$ ; à gauche : $q \leq p$}

\textbf{NB} Si on ne suppose pas que l'hom
\[ k\{T\} \to \mathrm{Gr}^*(S) \]
est un iso, on ne peut rien dire.
Prenons p.ex.\ le cas extrême et dégénéré où
\[ t = 0, \quad \text{donc } S \xrightarrow{\ \sim\ } k \]
on a alors
\[ \mathrm{Gr}^p\widehat{A}^0_* \simeq \Gamma^p_k(\Psi_*), \]
\struck{donc}
\[ \mathrm{Gr}^p(\widehat{A}^0_*) \simeq K_*(k, p) \]
qui n'est flasque pour aucun $p \geq 0$ (si $k \neq 0$).
Prenons aussi le cas de $(\mathbf{Z}_p, p)$, la
p.d.-filtration définie par $p$ est cofinale à la
filtration $p$-adique si $p \neq 2$ (si $p = 2$,
c'est la filtration constante de valeur
$\mathrm{Fil}^p\mathbf{Z}_2 = 2\mathbf{Z}_2$ si $p \geq 1$, donc $\mathrm{Gr}^*(S) \simeq \mathrm{Gr}^0(S) \simeq \mathbf{F}_2$
dans ce cas, i.e.\ c'est le cas de
dégénérescence…), quel est exactement l'idéal
$S^{[i]}$ engendré par les
\struck{$S^{[i]} = (\frac{p^n}{n!}\mathbf{Z}_p)$ -- i.e.\ quelle}
$\frac{p^j}{j!}$ pour $j \geq i$? \struck{i.e.\ probablement}
\note{le « $\simeq$ » de $\mathrm{Gr}^*(S) \simeq \mathrm{Gr}^0(S)$ est souligné ; la ligne biffée est écrite entre les deux dernières lignes de la question.}

\page{240}
\note{la page fait suite à la question qui clôt la p.\ 238 ; « $\mathrm{chiff}_p\,n$ » désigne, semble-t-il, la somme des chiffres de $n$ en base $p$.}
\[ v_p(n!) = \frac{1}{p-1}(n - \mathrm{chiff}_p\,n) \]
\[ v_p(p^{(n)}) = v_p\Bigl(\frac{p^n}{n!}\Bigr) = \frac{1}{p-1}\bigl((p-2)n + \underline{\mathrm{chiff}_p\,n}\bigr) \]
\[ \bigl(= \mathrm{chiff}_p(n) \text{ si } p = 2\bigr) \]
\struck{Si $p$}
\[ \gamma_{p,i} = \operatorname*{Inf}_{n \geq i} v_p(p^{(n)}) = \operatorname*{Inf}_{n \geq i} \frac{1}{p-1}\bigl(\underbrace{(p-2)n + \mathrm{chiff}_p(n)}_{\text{cette fonction n'est ni croissante ni décroissante…}}\bigr) \]
\[ = \begin{cases} 1 & \text{si } p = 2 \\ > \frac{p-2}{p-1}\,i \xrightarrow[i \to +\infty]{} +\infty & \text{si } p \neq 2 \end{cases} \]

\[ i = i_0 + i_1p + \dots + i_rp^r \]
\struck{donc} \struck{$i_0$} \struck{si}
\[ = \underbrace{\underbrace{(p-1)}_{i_0} + \dots + \underbrace{(p-1)p^{s-1}}_{i_{s-1}}}_{p^s - 1} + \underset{< p-1}{i_s}\,p^s + \dots + i_rp^r \]
\[ i + 1 = (i_s+1)p^s + \dots + i_rp^r \]
\[ (p-2)i + (p-1)s + \sum_{\alpha \geq s} i_\alpha \leq (p-2)(i+1) + \sum_{\alpha \geq s} i'_\alpha \]
\note{les deux sommes $\sum_{\alpha \geq s}$ sont barrées d'un trait oblique, comme pour les simplifier l'une par l'autre.}
\begin{align*}
[v_p(p^{(i)}) \leq v_p(p^{(i+1)})] &\Leftrightarrow [(p-1)s \leq (p-2)] \Leftrightarrow [s = 0] \Leftrightarrow [i_0 \neq p-1] \\
&\Leftrightarrow i + 1 \not\equiv 0 \ (p)
\end{align*}
\note{sous $[i_0 \neq p-1]$, une variante encadrée et biffée : $i_0 < p-1$.}

Exemple : si $\gamma_{p,i} = \gamma_{p,i+1}$: il \emph{suffit} que l'on
ait \struck{\ill{}} $v_p(p^{(i+1)}) \leq v_p(p^{(i)})$
i.e.\ que \struck{\ill{}} $i + 1 \equiv 0 \ (p)$
i.e.\ $i$ de la forme $p\ell - 1$.
\note{la page s'arrête ici ; l'étude des $\gamma_{p,i}$ se poursuit peut-être au-delà de ce lot.}

\end{document}
